% ==============================================
% ПРЕАМБУЛА ДЛЯ КОНСПЕКТОВ ПО ИНОСТРАННЫМ ЯЗЫКАМ
% ==============================================

% НАЧАЛО: Языки и шрифты
\usepackage{cmap}            % Поиск и копирование русского текста из PDF
\usepackage[utf8]{inputenc}  % Кодировка исходного текста
\usepackage[T1,T2A]{fontenc} % Поддержка кириллических шрифтов
\usepackage[romanian,russian]{babel} % Переносы строк по правилам языка
\usepackage{tipa}            % Фонетические знаки

% ГЕОМЕТРИЯ: Размеры полей страницы
\usepackage[
    left=2cm,
    right=2cm,
    top=2cm,
    bottom=2cm
]{geometry}

% ГРАФИКА: Иллюстрации и рисунки
\usepackage{graphicx}        % Вставка готовых картинок (PNG, JPG, PDF)
\usepackage{tikz}            % Рисование графиков и схем кодом

% ТИПОГРАФИКА: Оформление текста
\usepackage{indentfirst}     % Принудительный отступ (красная строка) для всех абзацев
\usepackage{microtype}       % Выравнивание краев текста без разрывов слов

% Междустрочный интервал
\usepackage{setspace}
\setstretch{1.0}

% Настройка абзацев
\usepackage[
    skip=0.5em,  % Вертикальный интервал (обычно 0.5em)
    indent=0em,  % Красная строка (обычно 1.5em)
]{parskip}

% Настройка списков
\usepackage{enumitem}

% Регистрируем русские буквы в пакете enumitem
\AddEnumerateCounter{\asbuk}{\@asbuk}{а}
\AddEnumerateCounter{\Asbuk}{\@Asbuk}{А}

\setlist{
    nosep,            % Убрать все вертикальные отступы
    leftmargin=1.5em, % Горизонтальный отступ от предыдущего уровня
}

% Маркеры списков (itemize)
\setlist[itemize]{label=---}

% Маркеры списков (enumerate)
\setlist[enumerate,1]{label=\arabic*.}
\setlist[enumerate,2]{label=\arabic*)}
\setlist[enumerate,3]{label=\asbuk*)}

% ЛОКАЛИЗАЦИЯ: Переименование стандартных названий
\addto\captionsrussian{\renewcommand{\chaptername}{Урок}}        % Глава
\addto\captionsrussian{\renewcommand{\contentsname}{Содержание}} % Оглавление

% НУМЕРАЦИЯ: Формат разделов
\setcounter{secnumdepth}{1} % Отключает нумерацию для subsection и ниже
\setcounter{tocdepth}{1}    % Убирает subsection и ниже из оглавления
\renewcommand{\thesection}{§\,\arabic{section}} % Номер раздела в виде параграфа: § 1

% ТАБЛИЦЫ
\usepackage{tabularray}    % Современные таблицы
\UseTblrLibrary{booktabs}  % Линии для таблиц

% Обнуление внешних отступов
\SetTblrOuter[tblr]{presep=0pt, postsep=0pt}
\SetTblrOuter[longtblr]{presep=0pt, postsep=0pt}

% Отключение подписей у длинных таблиц
\DefTblrTemplate{firsthead,middlehead,lasthead}{default}{}
\DefTblrTemplate{firstfoot,middlefoot,lastfoot}{default}{}

% КОНЕЦ: Интерактивность и ссылки
\usepackage[
    colorlinks=true,   % Цветной текст для ссылок вместо цветных рамок
    linkcolor=black,   % Черный цвет ссылок в оглавлении
    urlcolor=blue,     % Синий цвет для интернет-ссылок
    citecolor=black    % Черный цвет для библиографии
]{hyperref}
